\section{Esportazione}
\label{sec:esportazione}

La mesh deve essere esportata in formato GLTF. Per questo compito è stata utilizzata la libreria tinygltf \cite{tinygltf}. Vengono esportate di base le informazioni sui vertici: posizione, normale e coordinate texture.

Attualmente, porte e finestre sono rappresentate come semplici aperture nei muri: questo, tuttavia, non è sufficiente, poiché si vuole poter inserire, in corrispondenza di tali aperture, degli asset 3D di porte e finestre.

Una prima ipotesi era inserire gli asset direttamente nel processo di generazione della mesh, ma questo approccio risulta poco conveniente: se in un appartamento sono presenti dieci porte, la geometria della porta verrebbe duplicata altrettante volte, aumentando notevolmente la dimensione del file, tanto più se si considera che un modello dettagliato può contare centinaia di migliaia di triangoli. Inoltre, se si vuole cambiare asset o modificarne l'orientamento, sarebbe necessario rieseguire l'intera pipeline. Infine, occorrerebbe anche preservare i materiali originali dell'asset scaricato, complicando ulteriormente il processo.

La soluzione più flessibile e mantenibile consiste nell'adottare un approccio a placeholder, spostando il caricamento degli asset direttamente in Blender. L'idea è di esportare le informazioni sui varchi: quando una faccia viene riconosciuta come porta o finestra, se ne calcolano posizione, larghezza, altezza e spessore, oltre alla rotazione. Più precisamente, il centro rappresenta il centro dell'apertura, calcolato tenendo conto di stipite e davanzale: non si trova quindi all'altezza del soffitto, bensì esattamente a metà dell'apertura stessa.

Tutte queste informazioni vengono infine serializzate ed esportate sotto forma di file JSON. Ad esempio:
\begin{listing}[H]
\begin{minted}[frame=single]{json}
{
	"openings": [
	{
		"type": "DOOR",
		"center": [12.5, 0.0, 3.2],
		"width": 0.9,
		"height": 2.1,
		"thickness": 0.1,
		"rotation_z": 90.0
	},
	{
		"type": "WINDOW",
		"center": [8.3, 0.0, 1.5],
		"width": 1.2,
		"height": 1.4,
		"thickness": 0.1,
		"rotation_z": 0.0
	}]
}
\end{minted}
\end{listing}